\section{Separating existence, detectability and observation}
\label{sec:decomp}

\subsection{The decomposition}

Let $\Ntrue$ be the number of technological civilizations that exist in the
relevant volume and epoch. Write $\Psurv$ for the probability that such a
civilization persists into the phase we could in principle detect, $h(\Obs)$ for
the probability that it produces a signature strong enough to register given its
external observability $\Obs$, and $\Psearch$ for the probability that our
searches intersect that signature. Then
\begin{equation}
  \boxed{\;\Nobs \;=\; \Ntrue \,\cdot\, \Psurv \,\cdot\, h(\Obs) \,\cdot\, \Psearch\;}
  \label{eq:decomp}
\end{equation}
with $h$ a saturating map from signature strength to detection probability. We
take
\begin{equation}
  h(\Obs) \;=\; 1 - e^{-\kappa \Obs},
  \label{eq:hdet}
\end{equation}
where $\kappa$ sets instrument sensitivity, and define the detection probability
for a single civilization as $\Pdet = h(\Obs)\,\Psearch$.

Equation~\eqref{eq:decomp} is a bookkeeping identity, not a theory; its value is
that it forces the four mechanisms apart. \Established{} Figure~\ref{fig:venn}
shows the resulting nesting and Figure~\ref{fig:drakemap} shows where the
framework acts relative to Eq.~\eqref{eq:drake}.

\begin{figure}[t]
\centering
\begin{tikzpicture}[font=\small]
  \draw[rounded corners=8pt, fill=blue!5, draw=blue!50]
        (0,0) rectangle (10.4,2.6);
  \node[anchor=north west] at (0.15,2.55) {$\Ntrue$: civilizations that exist};
  \draw[rounded corners=6pt, fill=blue!12, draw=blue!60]
        (0.7,0.35) rectangle (9.7,1.95);
  \node[anchor=north west] at (0.85,1.9) {$\Ndet$: produce a detectable signature ($\times\,\Psurv\, h(\Obs)$)};
  \draw[rounded corners=4pt, fill=blue!25, draw=blue!70]
        (1.5,0.6) rectangle (6.0,1.35);
  \node[anchor=west] at (1.65,0.98) {$\Nobs$: actually observed ($\times\,\Psearch$)};
  \node[anchor=west, font=\itshape\footnotesize] at (6.3,0.98) {non-observation lives here};
  \draw[-{Latex[length=2mm]}, thick] (6.05,0.98) -- (6.25,0.98);
\end{tikzpicture}
\caption{Existence, detectability and observation are distinct. Non-observation
constrains the innermost set, and therefore only the \emph{product} of the
factors in Eq.~\eqref{eq:decomp}; it does not by itself constrain $\Ntrue$.}
\label{fig:venn}
\end{figure}

\begin{figure}[t]
\centering
\begin{tikzpicture}[font=\small, node distance=3mm]
  \node (drake) [draw, rounded corners, fill=gray!8, minimum height=8mm]
    {$N = R_*\, f_p\, n_e\, f_l\, f_i\, f_c\, L$};
  \node (lab1) [below=1mm of drake, font=\footnotesize\itshape] {classical Drake equation};
  \node (arrow) [below=7mm of lab1] {};
  \node (rof) [draw, rounded corners, fill=blue!8, minimum height=8mm, below=9mm of lab1]
    {$\Nobs = \Ntrue \cdot \Psurv \cdot h(\Obs) \cdot \Psearch$};
  \node (lab2) [below=1mm of rof, font=\footnotesize\itshape] {observability-aware form};
  \draw[-{Latex[length=2mm]}, thick] (lab1) -- node[right, font=\footnotesize] {refine $L$} (rof);
  \node (act) [right=14mm of rof, draw, dashed, rounded corners, fill=orange!10,
               text width=3.3cm, align=center, font=\footnotesize]
    {the model acts here:\\ $\Psurv$ and $h(\Obs)$};
  \draw[-{Latex[length=2mm]}, dashed] (act) -- (rof);
\end{tikzpicture}
\caption{The framework is a refinement of the Drake equation's lifetime and
detectability content, not a replacement for it. It leaves the astrophysical
factors untouched and acts on survival and signal production.}
\label{fig:drakemap}
\end{figure}

\subsection{Why this matters for inference}

Suppose a survey reports no detections. Under Eq.~\eqref{eq:decomp} the
likelihood of that outcome depends on $\Psearch$ and on the distribution of
$\Obs$ across the population, not on $\Ntrue$ alone. Two very different worlds
--- one in which civilizations are rare, and one in which they are common but
mostly quiet during the epoch we can sample --- produce the same null result.
Distinguishing them requires a model of how $\Obs$ is distributed and how it
evolves, which is what \S\ref{sec:model} supplies. \Established

\subsection{What observability is, and is not}

We use $\Obs$ for \emph{external observability}: the strength of signatures
leaving the system, aggregating waste heat, leakage and intentional emission,
expansion footprint, and engineered structures. It is not energy use, and it is
not capability. A civilization may be extremely capable and comparatively quiet;
it may also be modest and loud.

Crucially, $\Obs$ can be small but never zero. Any energy use produces
thermodynamic consequences, and a claim of true invisibility would require
physics we have no warrant to assume \citep{dyson1960,kardashev1964}. We encode
this directly: $\Obs$ is bounded below by a strictly positive floor $\Ofloor$
(Proposition~\ref{prop:obs}). Low observability in this framework always means
``hard to detect'', never ``absent''. \Established
